\documentclass[UTF8,zihao=-4]{ctexart}
\usepackage[a4paper,margin=1.75cm]{geometry}
\usepackage{amsmath,amssymb}
\usepackage{booktabs}
\usepackage{longtable}
\usepackage{array}
\usepackage{xcolor}
\usepackage{hyperref}

\hypersetup{colorlinks=true,linkcolor=blue!60!black,urlcolor=blue!60!black}
\setlength{\parindent}{0pt}
\setlength{\parskip}{0.45em}
\renewcommand{\arraystretch}{1.22}

\newcommand{\ans}[1]{\textcolor{blue!60!black}{#1}}
\newcommand{\warn}[1]{\textcolor{red!70!black}{#1}}

\title{\bfseries 第5章理想流体：原题练习汇总答案版}
\author{}
\date{}

\begin{document}
\maketitle

本答案版对应《第5章理想流体：原题练习汇总（去作业版）》。
课本中有些题是建模题或证明题，答案给出关键方程、边界条件或结论；数值题给最终结果。

\section*{一、先记住的基本解}

\begin{longtable}{p{3cm}p{5.1cm}p{5.8cm}}
\toprule
\textbf{基本流动} & \textbf{速度势 $\varphi$} & \textbf{流函数 $\psi$}\\
\midrule
均匀流 $U$ 沿 $+x$ & $\varphi=Ux$ & $\psi=Uy$\\
点源 $q$ 在原点 & $\displaystyle \varphi=\frac{q}{2\pi}\ln r$ & $\displaystyle \psi=\frac{q}{2\pi}\theta$\\
点涡 $\Gamma$ 在原点 & $\displaystyle \varphi=\frac{\Gamma}{2\pi}\theta$ & $\displaystyle \psi=-\frac{\Gamma}{2\pi}\ln r$\\
偶极子 & 按题中方向由点源、点汇极限得到 & 按题中方向由点源、点汇极限得到\\
\bottomrule
\end{longtable}

总原则：
\[
\psi=\psi_1+\psi_2+\cdots,\qquad
\varphi=\varphi_1+\varphi_2+\cdots .
\]
流线方程：
\[
\psi=C.
\]
驻点条件：
\[
u=0,\qquad v=0.
\]

\section*{二、课本第5章习题参考答案}

\begin{longtable}{p{1.3cm}p{12.6cm}}
\toprule
\textbf{题号} & \textbf{答案/关键结论}\\
\midrule
\endfirsthead
\toprule
\textbf{题号} & \textbf{答案/关键结论}\\
\midrule
\endhead

5.1 &
速度势提法：
\[
\nabla^2\varphi=0.
\]
边界条件：物面不可穿透，
\[
\frac{\partial\varphi}{\partial n}=\vec V_{\text{物面}}\cdot\vec n;
\]
无穷远处满足给定来流。题设 $\Gamma=0$，所以不再附加环量条件。\\

5.2 &
流函数提法：
\[
\nabla^2\psi=0.
\]
壁面为流线，故壁面上 $\psi=\text{常数}$；进口、出口的流函数差等于体积流量 $Q$。\\

5.3 &
绝对流动无旋，所以仍满足势流/流函数的拉普拉斯方程。边界条件要写在随椭圆柱运动的坐标系中：物面法向相对速度为零，无穷远处为与物体运动相反的来流。\\

5.4 &
若题面为
\[
u=x^2+y^2,\qquad v=x^2-2xy,
\]
则
\[
\frac{\partial u}{\partial x}+\frac{\partial v}{\partial y}=2x-2x=0,
\]
所以存在不可压流函数；但
\[
\frac{\partial v}{\partial x}-\frac{\partial u}{\partial y}
=2x-4y\ne0,
\]
所以不存在速度势。

流函数可取：
\[
\psi=x^2y+\frac{y^3-x^3}{3}.
\]

变形速率：
\[
e_{xx}=2x,\qquad e_{yy}=-2x,\qquad e_{xy}=x.
\]
刚体旋转角速度：
\[
\Omega_z=\frac12\left(\frac{\partial v}{\partial x}-\frac{\partial u}{\partial y}\right)=x-2y.
\]\\

5.5 &
总流函数：
\[
\psi=Uy+\frac{Q}{2\pi}\theta_1+\frac{Q}{2\pi}\theta_2.
\]
若 $U=5,\ Q=20\pi$，驻点在
\[
(x,y)=(-2,0)
\]
附近；过驻点流线可写为
\[
5y+10(\theta_1+\theta_2)=0.
\]
具体方向以原题图中来流方向为准；若来流反向，驻点关于 $y$ 轴对称变号。\\

5.6 &
见第三节完整写法。核心结论：
\[
U_\infty=\frac{\Gamma}{4\pi h},
\]
\[
\ln\frac{\sqrt{x^2+(y+h)^2}}{\sqrt{x^2+(y-h)^2}}
-\frac{y}{2h}=C.
\]\\

5.7 &
用均匀流、点源和连续分布点汇叠加。写出总流函数 $\psi$ 后，物面由过驻点流线确定：
\[
\psi=C_{\text{驻点}}.
\]
求法：先由 $u=0,\ v=0$ 找驻点，再令 $\psi=\psi_{\text{驻点}}$。\\

5.8 &
书后答案给出关键结论：常数 $c$ 由无穷远来流和奇点强度匹配确定；最大速度点在对称轴附近。该题重点不是套公式，而是用轴对称基本解叠加。\\

5.9 &
这是兰金卵形体绕流。物面方程由
\[
\psi=\psi_{\text{驻点}}
\]
确定；最大速度点位于对称轴上。\\

5.10 &
若
\[
\psi=\arctan\frac{y}{x},
\]
则
\[
w(z)=\ln z.
\]
若
\[
\psi=\ln(x^2+y^2),
\]
则
\[
w(z)=2i\ln z.
\]\\

5.11 &
若
\[
\varphi=r^2\cos2\theta,
\]
则
\[
w(z)=z^2.
\]
第二小问按圆柱绕流型势函数写复势，常见形式为
\[
w(z)=Ue^{-i\alpha}\left(z+\frac{a^2}{z}\right),
\]
符号按题面来流方向调整。\\

5.12 &
按图中点源、点汇写复势并相加：
\[
w(z)=\sum_k \frac{q_k}{2\pi}\ln(z-z_k).
\]
证明 $|z|=a$ 是流线，只需证明在 $|z|=a$ 上
\[
\operatorname{Im}w(z)=\text{常数}.
\]\\

5.13 &
点涡轨迹等于它所在流场的流线。代入初始点 $(\sqrt2,\sqrt2)$ 确定常数即可。\\

5.14 &
书后结果：
\[
Q=12\pi,\qquad \Gamma=8\pi.
\]\\

5.15 &
用保角映射把圆弧边界化为较简单边界，再写复势。此题主要练映射，不是普通代数题。\\

5.16 &
书后结果：
\[
\Gamma=-20.5\ \mathrm{m^2/s},
\]
\[
x_{\text{驻点}}\approx -0.74\ \mathrm{m},
\]
\[
|F|\approx125\ \mathrm{N},
\qquad F_x\approx -10.9\ \mathrm{N},\quad F_y\approx124.5\ \mathrm{N}.
\]\\

5.17 &
无环量椭圆柱势流绕流：
\[
F_x=0,\qquad F_y=0,
\]
但有力矩：
\[
M_z=-\rho U^2\pi(a^2-b^2)\sin2\alpha.
\]
符号按原题坐标正向调整。\\

5.18 &
可压缩平面定常流动的连续方程：
\[
\frac{\partial(\rho u)}{\partial x}
+\frac{\partial(\rho v)}{\partial y}=0.
\]
因此可定义流函数：
\[
\rho u=\frac{\partial\psi}{\partial y},\qquad
\rho v=-\frac{\partial\psi}{\partial x}.
\]\\

5.19 &
按极坐标写流线方程 $\theta=\theta(r)$，速度方向与流线切向一致，再由
\[
\omega_z=\frac1r\frac{\partial(rv_\theta)}{\partial r}
-\frac1r\frac{\partial v_r}{\partial\theta}
\]
推出题中涡量表达式。\\

5.20 &
三个等强点涡位于等边三角形三个顶点时，三角形形状保持不变，整体绕面心匀速旋转。
若外接圆半径为 $R$，点涡强度为 $\Gamma$，则角速度量级为
\[
\Omega=\frac{\Gamma}{2\pi R^2}.
\]\\

5.21 &
中心点源只产生径向速度，不能抵消三个同号点涡相互诱导产生的切向旋转速度；若题面确为“点源”，则不存在有限 $Q$ 使三点涡完全静止。若原题实际为中心点涡，则需用中心涡诱导切向速度抵消 5.20 的旋转速度。\\

5.22 &
用不可压平面理想有旋流动的涡量输运方程：
\[
\frac{D\omega}{Dt}=0
\]
和随体面积元不变，直接对题中积分作质点导数，可得
\[
\frac{D I_m}{Dt}=0.
\]\\

S5.1 &
会旋转。平板下落时微小倾斜会造成两侧压力和速度分布不对称，从而产生力矩。\\

S5.2 &
翼尖处上下翼面压力差使气流绕过翼尖，形成一对尾涡；两条涡管旋转方向相反。起动时后缘脱落起动涡，由环量守恒，绕翼型产生相反方向环量。\\

S5.3 &
奇点必须放在物体内部，否则会在流体区域引入非物理奇点；总点源强度等于总点汇强度，是为了保证远处没有净流量增减。圆外点源关于圆的镜像必须配成源汇，是为了同时满足圆面为流线和无穷远条件。\\

S5.4 &
不可压平面理想流动中
\[
\frac{D\omega}{Dt}=0.
\]
定常时质点轨迹与流线重合，所以涡量沿流线为常数。轴对称有旋流中不再是 $\omega$ 沿流线常数，而是
\[
\frac{\omega}{r}=\text{沿流线常数}.
\]\\

S5.5 &
对称翼型零攻角上下对称，环量为零，所以升力为零。有攻角时由库塔条件确定非零环量，故有升力。起飞低速时为获得足够升力，攻角应大于高空平飞时攻角。\\

\bottomrule
\end{longtable}

\section*{三、5.6 完整答案}

题意：两个等强反向点涡放在 $(0,h)$ 和 $(0,-h)$，远处来流刚好使两个涡停留不动。

上方点涡受到下方点涡的诱导速度：
\[
u_{\text{ind}}=\frac{\Gamma}{2\pi(2h)}
=\frac{\Gamma}{4\pi h}.
\]
来流必须与此相反，大小为
\[
U_\infty=\frac{\Gamma}{4\pi h}.
\]

取上方点涡为 $+\Gamma$，下方点涡为 $-\Gamma$，且均匀来流向左，则
\[
\psi_\infty=-U_\infty y.
\]
两点涡的流函数为
\[
\psi_1=-\frac{\Gamma}{2\pi}\ln\sqrt{x^2+(y-h)^2},
\]
\[
\psi_2=+\frac{\Gamma}{2\pi}\ln\sqrt{x^2+(y+h)^2}.
\]
总流函数：
\[
\psi
=-U_\infty y
-\frac{\Gamma}{2\pi}\ln\sqrt{x^2+(y-h)^2}
+\frac{\Gamma}{2\pi}\ln\sqrt{x^2+(y+h)^2}.
\]
代入 $U_\infty=\Gamma/(4\pi h)$：
\[
\psi=
\frac{\Gamma}{2\pi}
\left[
\ln
\frac{\sqrt{x^2+(y+h)^2}}
{\sqrt{x^2+(y-h)^2}}
-\frac{y}{2h}
\right].
\]
所以流线方程为
\[
\boxed{
\ln
\frac{\sqrt{x^2+(y+h)^2}}
{\sqrt{x^2+(y-h)^2}}
-\frac{y}{2h}=C
}.
\]

\section*{四、往年题答案}

\subsection*{1. 不可压缩无旋线性速度场}

题面：
\[
u=Ax+By,\qquad v=Cx+Dy,\qquad w=0.
\]

不可压缩连续方程：
\[
\frac{\partial u}{\partial x}+\frac{\partial v}{\partial y}=A+D=0.
\]

无旋条件：
\[
\frac{\partial v}{\partial x}-\frac{\partial u}{\partial y}=C-B=0.
\]

所以
\[
\boxed{D=-A,\qquad C=B.}
\]

\subsection*{2. 2022 春期末：半圆柱气膜馆势流}

把地面看作对称面，半圆柱外流等价于完整圆柱绕流的上半部分。

圆柱表面速度：
\[
v_r=0,\qquad v_\theta=-2V\sin\theta.
\]

由伯努利：
\[
p(\theta)+\frac12\rho v_\theta^2
=p_\infty+\frac12\rho V^2.
\]
所以表面压力分布：
\[
\boxed{
p(\theta)=p_\infty+\frac12\rho V^2(1-4\sin^2\theta)
}.
\]

升力按单位长度计算，取动压相对压力积分：
\[
L'=-a\int_0^\pi [p(\theta)-p_\infty]\sin\theta\,d\theta.
\]
代入压力分布：
\[
L'
=-a\frac12\rho V^2
\int_0^\pi (1-4\sin^2\theta)\sin\theta\,d\theta
=\frac53\rho V^2a.
\]

代入
\[
a=5\ \mathrm{m},\qquad V=10\ \mathrm{m/s},\qquad \rho=1.29\ \mathrm{kg/m^3},
\]
得
\[
\boxed{
L'\approx1.08\times10^3\ \mathrm{N/m}
}
\]
方向向上。

\end{document}
